% !TeX program = xelatex
% Encoding: UTF-8
\documentclass[12pt,a4paper]{article}

\usepackage{fontspec}
\usepackage{amsmath,amssymb,amsthm,mathtools}
\usepackage{geometry}
\usepackage{enumitem}
\usepackage{hyperref}

\geometry{margin=2.0cm}
\setmainfont{Latin Modern Roman}

\setlength{\parindent}{2em}
\setlength{\parskip}{0.5em}

% Make displayed matrices easier to read
\renewcommand{\arraystretch}{1.35}
\setlength{\arraycolsep}{9pt}

\newtheorem{theorem}{Theorem}

\title{\textbf{Schur's Theorem}\\
\large Proof by Mathematical Induction}
\author{}
\date{}

\begin{document}
\maketitle

\section{Unitary Matrix}

For an $n\times n$ complex matrix $Q$, if
\[
    \overline{Q}^{\,T}=Q^{-1},
\]
then $Q$ is called a \textbf{unitary matrix}.

If a unitary matrix is real, then
\[
    \overline{Q}=Q,
\]
and hence
\[
    Q^T=Q^{-1}.
\]
Thus, in the real case, a unitary matrix is an \textbf{orthogonal matrix}.

\section{Schur's Theorem}

\begin{theorem}[Schur's Theorem]
    Every square matrix can be factored as
    \[
        \boxed{
            A=QTQ^{-1}=QT\overline{Q}^{\,T},
        }
    \]
    where $T$ is upper triangular and $Q$ is unitary.

    If $A$ has real eigenvalues, then $Q$ and $T$ can be chosen real. In this case,
    \[
        Q^T=Q^{-1},
    \]
    so $Q$ is orthogonal.
\end{theorem}

\section{Proof}

We prove the theorem by mathematical induction.

\subsection*{1. Base Case}

For $n=1$,
\[
    a=1\cdot a\cdot 1^{-1}.
\]

\subsection*{2. Induction Hypothesis}

Assume the theorem holds for every $k\times k$ matrix, and let $A$ be a
$(k+1)\times(k+1)$ matrix.

Let $\lambda_1$ be an eigenvalue of $A$ and $q_1$ a corresponding unit eigenvector:
\[
    Aq_1=\lambda_1q_1,
    \qquad
    \|q_1\|=1.
\]

Using Gram--Schmidt, extend $q_1$ to an orthonormal basis
\[
    q_1,q_2,\ldots,q_{k+1}.
\]

Let
\[
    Q_1=
    \left[
        \begin{array}{cccc}
            q_1 & q_2 & \cdots & q_{k+1}
        \end{array}
        \right].
\]
Then $Q_1$ is unitary.

\subsection*{3. Compute $\overline{Q}_1^{\,T}AQ_1$}

Write the multiplication explicitly:
\[
    \overline{Q}_1^{\,T}AQ_1
    =
    \left[
        \begin{array}{c}
            \overline{q}_1^{\,T} \\[3pt]
            \overline{q}_2^{\,T} \\[3pt]
            \vdots               \\[3pt]
            \overline{q}_{k+1}^{\,T}
        \end{array}
        \right]
    \left[
        \begin{array}{cccc}
            Aq_1 & Aq_2 & \cdots & Aq_{k+1}
        \end{array}
        \right].
\]

Since
\[
    Aq_1=\lambda_1q_1,
\]
we get
\[
    \overline{Q}_1^{\,T}AQ_1
    =
    \left[
        \begin{array}{cccc}
            \overline{q}_1^{\,T}\lambda_1q_1     &
            \overline{q}_1^{\,T}Aq_2             &
            \cdots                               &
            \overline{q}_1^{\,T}Aq_{k+1}
            \\[5pt]
            \overline{q}_2^{\,T}\lambda_1q_1     &
            \overline{q}_2^{\,T}Aq_2             &
            \cdots                               &
            \overline{q}_2^{\,T}Aq_{k+1}
            \\[5pt]
            \vdots                               & \vdots & \ddots & \vdots
            \\[5pt]
            \overline{q}_{k+1}^{\,T}\lambda_1q_1 &
            \overline{q}_{k+1}^{\,T}Aq_2         &
            \cdots                               &
            \overline{q}_{k+1}^{\,T}Aq_{k+1}
        \end{array}
        \right].
\]

Because the $q_i$ are orthonormal,
\[
    \overline{q}_1^{\,T}q_1=1,
    \qquad
    \overline{q}_j^{\,T}q_1=0
    \quad(j=2,\ldots,k+1).
\]

Hence the first column becomes
\[
    \left[
        \begin{array}{c}
            \lambda_1 \\[4pt]
            0         \\[4pt]
            \vdots    \\[4pt]
            0
        \end{array}
        \right],
\]
and therefore
\[
    \boxed{
        \overline{Q}_1^{\,T}AQ_1
        =
        \left[
            \begin{array}{ccccc}
                \lambda_1 & * & *   & \cdots & * \\[5pt]
                0         &   &     &        &   \\[-3pt]
                0         &   & A_2 &        &   \\[-3pt]
                \vdots    &   &     &        &   \\[-3pt]
                0         &   &     &        &
            \end{array}
            \right].
    }
\]

Here $A_2$ is a $k\times k$ matrix.

\subsection*{4. Apply the Induction Hypothesis to $A_2$}

By the induction hypothesis,
\[
    A_2=Q_2T_2\overline{Q}_2^{\,T},
\]
where $Q_2$ is unitary and $T_2$ is upper triangular. Hence
\[
    \overline{Q}_2^{\,T}A_2Q_2=T_2.
\]

Let
\[
    Q=
    Q_1
    \left[
        \begin{array}{ccccc}
            1      & 0 & 0   & \cdots & 0 \\
            0      &   &     &        &   \\
            0      &   & Q_2 &        &   \\
            \vdots &   &     &        &   \\
            0      &   &     &        &
        \end{array}
        \right].
\]

Denote the second matrix by $D$, so that $Q=Q_1D$.

\subsection*{5. Show that $Q$ is Unitary}

Since $Q_2$ is unitary,
\[
    Q_2\overline{Q}_2^{\,T}=I_k,
\]
and hence
\[
    D\overline{D}^{\,T}=I_{k+1}.
\]

Also,
\[
    Q_1\overline{Q}_1^{\,T}=I_{k+1}.
\]

Therefore
\[
    Q\overline{Q}^{\,T}
    =
    Q_1D\overline{D}^{\,T}\overline{Q}_1^{\,T}
    =
    I_{k+1}.
\]

Thus $Q$ is unitary.

\subsection*{6. Compute $\overline{Q}^{\,T}AQ$}

We have
\[
    \overline{Q}^{\,T}AQ
    =
    \overline{D}^{\,T}
    \overline{Q}_1^{\,T}AQ_1
    D.
\]

Writing out the matrices,
\[
    \overline{Q}^{\,T}AQ
    =
    \left[
        \begin{array}{ccccc}
            1      & 0 & 0                    & \cdots & 0 \\
            0      &   &                      &        &   \\
            0      &   & \overline{Q}_2^{\,T} &        &   \\
            \vdots &   &                      &        &   \\
            0      &   &                      &        &
        \end{array}
        \right]
    \left[
        \begin{array}{ccccc}
            \lambda_1 & * & *   & \cdots & * \\
            0         &   &     &        &   \\
            0         &   & A_2 &        &   \\
            \vdots    &   &     &        &   \\
            0         &   &     &        &
        \end{array}
        \right]
    \left[
        \begin{array}{ccccc}
            1      & 0 & 0   & \cdots & 0 \\
            0      &   &     &        &   \\
            0      &   & Q_2 &        &   \\
            \vdots &   &     &        &   \\
            0      &   &     &        &
        \end{array}
        \right].
\]

Therefore
\[
    \overline{Q}^{\,T}AQ
    =
    \left[
        \begin{array}{ccccc}
            \lambda_1 & * & *                          & \cdots & * \\
            0         &   &                            &        &   \\
            0         &   & \overline{Q}_2^{\,T}A_2Q_2 &        &   \\
            \vdots    &   &                            &        &   \\
            0         &   &                            &        &
        \end{array}
        \right].
\]

Since
\[
    \overline{Q}_2^{\,T}A_2Q_2=T_2,
\]
we obtain
\[
    \boxed{
        \overline{Q}^{\,T}AQ
        =
        \left[
            \begin{array}{ccccc}
                \lambda_1 & * & *   & \cdots & * \\
                0         &   &     &        &   \\
                0         &   & T_2 &        &   \\
                \vdots    &   &     &        &   \\
                0         &   &     &        &
            \end{array}
            \right]
        =T.
    }
\]

Since $T_2$ is upper triangular, $T$ is upper triangular. Thus
\[
    \overline{Q}^{\,T}AQ=T.
\]

Because
\[
    \overline{Q}^{\,T}=Q^{-1},
\]
we have
\[
    Q^{-1}AQ=T,
\]
and therefore
\[
    \boxed{
        A=QTQ^{-1}=QT\overline{Q}^{\,T}.
    }
\]

Hence Schur's theorem is proved.
\end{document}
